% Memo -- example for the plaindoc class. People and products are fictional.
% Build: latexmk -lualatex datasheet-change-memo.tex
\documentclass[type=memo, lang=en-GB]{plaindoc}

\pdset{
  org   = {Orrin Instruments (fictional)},
  title = {Derate DP305 current above 35\,°C},
  to    = {Product Marketing},
  from  = {M. Chan, Product Engineering},
  cc    = {Quality, Technical Documentation},
  date  = {25 November 2026},
  docid = {TR-2026-041},
}

\begin{document}
\maketitle

Please change the DP305 datasheet and user manual before the January print
run. Thermal testing (TR-2026-041) shows that the DP305 cannot deliver its
full \qty{5}{\ampere} for longer than about an hour at \qty{40}{\degreeCelsius}
ambient. Below \qty{35}{\degreeCelsius} it is unaffected.

\section*{What changes}

\begin{itemize}
\item Specifications: add \enquote{Full output up to \qty{35}{\degreeCelsius}
      ambient; derate current linearly to \qty{4}{\ampere} at
      \qty{40}{\degreeCelsius}.}
\item User manual, chapter 7: the same sentence under \emph{Operating temperature}.
\item No firmware or hardware change. Over-temperature protection already
      switches the output off safely.
\end{itemize}

\section*{Dates}

\begin{tabularx}{\linewidth}{@{}lL@{}}
\toprule
\hd{Date} & \hd{Action} \\
\midrule
5 December 2026  & Engineering confirms the \qty{4}{\ampere} limit on three units. \\
12 December 2026 & Marketing and Documentation submit the revised pages for review. \\
9 January 2027   & Print run with the revised datasheet. \\
\bottomrule
\end{tabularx}

\medskip
Units already shipped do not need any action. The full test report is
attached.

\end{document}
